%%% SDGFT-2026-02
%%% Topology of the 24-Cell and F4 Symmetry in Scale-Dependent Geometric
%%% Field Theory
%%%
\documentclass[amsmath,amssymb,aps,prd,twocolumn,superscriptaddress,nofootinbib]{revtex4-2}
\usepackage{graphicx}
\usepackage{hyperref}
\usepackage{xcolor}
\usepackage{booktabs}
\usepackage{float}
\usepackage{sdgft-macros}

\begin{document}

\preprint{SDGFT-2026-02}

\title{Topology of the 24-Cell and $F_4$ Symmetry in Scale-Dependent
       Geometric Field Theory}

\author{David~A.~Besemer}
\email{david.besemer@sdgft.org}
\affiliation{Independent Researcher}
\date{\today}

\begin{abstract}
We investigate the topology and symmetry of the 24-cell polytope
$\{3,4,3\}$, the fundamental geometric object of Scale-Dependent Geometric
Field Theory (SDGFT).  After reviewing the vertex coordinates on the unit
three-sphere $S^3$ and establishing the self-duality of the polytope, we
decompose the 24-cell into three mutually orthogonal 16-cells, each
contributing eight vertices ($8+8+8=24$).  The full automorphism group is the
Weyl group $W(F_4)$ of order~1152.  We discuss $D_4$~triality---the
order-three outer automorphism of $\mathfrak{so}(8)$---and show how it
permutes the three 16-cell sub-polytopes.  Finally, we trace the Dynkin-diagram
chain $F_4 \to D_4 \to A_2 \oplus A_1 \oplus U(1)$, recovering the Standard
Model gauge group $SU(3)_C \times SU(2)_L \times U(1)_Y$ via successive
peripheral-node removal.
\end{abstract}

\maketitle

%======================================================================
\section{Introduction}\label{sec:intro}
%======================================================================

The 24-cell occupies a privileged position in four-dimensional geometry: it is
the unique self-dual regular polytope with no lower-dimensional counterpart.
In Scale-Dependent Geometric Field Theory~\cite{Besemer2026_01}, every physical
constant is derived from the combinatorial and topological data of this single
object.  The present paper provides the mathematical underpinning for that
programme by analysing the 24-cell's geometry, symmetry group, and
representation-theoretic structure in detail.

Section~\ref{sec:geometry} reviews the geometric properties of the 24-cell,
including its vertex coordinates and self-duality proof.
Section~\ref{sec:decomp} presents the decomposition into three orthogonal
16-cells.  Sections~\ref{sec:F4}--\ref{sec:dynkin} develop the $F_4$ symmetry,
$D_4$ triality, and the Dynkin-diagram descent to the Standard Model gauge
group.  Section~\ref{sec:conclusions} concludes.

%======================================================================
\section{Geometric Properties}\label{sec:geometry}
%======================================================================

\subsection{Vertex coordinates on $S^3$}

The 24 vertices of the 24-cell may be placed on the unit three-sphere
$S^3 \subset \mathbb{R}^4$.  They fall into two orbits under the hyperoctahedral
group $B_4$:
\begin{align}
  V_1 &= \bigl\{(\pm1,0,0,0)\text{ and permutations}\bigr\}
         \quad(8\text{ vertices})\,,\label{eq:V1}\\[4pt]
  V_2 &= \bigl\{\tfrac{1}{\sqrt{2}}(\pm1,\pm1,0,0)
         \text{ and permutations}\bigr\}
         \quad(16\text{ vertices})\,.\label{eq:V2}
\end{align}
Every vertex has unit norm, $|v|=1$, and the minimal edge length is
$\ell = \sqrt{2 - \sqrt{2}} \approx 1$ after rescaling.

\subsection{Face and cell counts}

The $f$-vector of the 24-cell is
\begin{equation}\label{eq:fvec}
  (f_0,\,f_1,\,f_2,\,f_3) = (24,\,96,\,96,\,24)\,.
\end{equation}
The equality $f_0 = f_3$ and $f_1 = f_2$ is the hallmark of a self-dual
polytope.

\subsection{Self-duality proof}

A polytope $P$ is \emph{self-dual} if there exists an isomorphism
$\phi\colon P \to P^*$ between $P$ and its dual $P^*$ that reverses the
face lattice.  For the 24-cell, the dual is constructed by placing a vertex at
the centroid of each octahedral cell.  Because every octahedral cell of
$\{3,4,3\}$ is regular, its centroid lies at a well-defined point on $S^3$.
The 24~centroids form a polytope combinatorially isomorphic to the original,
with $f$-vector $(24,96,96,24)$.  This establishes
\begin{equation}\label{eq:selfdual}
  \{3,4,3\}^* \cong \{3,4,3\}\,.
\end{equation}

More precisely, the dual vertices are obtained by the inversion map
$v \mapsto v/|v|^2$ (which acts as the identity on $S^3$), followed by a
rotation in $\mathrm{SO}(4)$.  The existence of this rotation is guaranteed by
the transitive action of $W(F_4)$ on cells and vertices alike~\cite{Coxeter1973}.

The self-duality of the 24-cell is unique among regular 4-polytopes: the
5-cell is self-dual but exists in all dimensions (as the simplex), while the
hypercube $\{4,3,3\}$ and its dual 16-cell $\{3,3,4\}$ are \emph{not}
self-dual.  The 120-cell and 600-cell form a dual pair but are not individually
self-dual.

%======================================================================
\section{Decomposition into Three 16-Cells}\label{sec:decomp}
%======================================================================

The 24 vertices of the 24-cell admit a partition into three mutually orthogonal
16-cells (cross-polytopes $\beta_4$).  Concretely, label the three subsets as
\begin{align}
  C_1 &: \{(\pm1,0,0,0),\;(0,\pm1,0,0),\;(0,0,\pm1,0),\;(0,0,0,\pm1)\}\,,
         \label{eq:C1}\\[3pt]
  C_2 &: \tfrac{1}{\sqrt{2}}\{(\pm1,\pm1,0,0),\;(0,0,\pm1,\pm1)\}\,,
         \label{eq:C2}\\[3pt]
  C_3 &: \tfrac{1}{\sqrt{2}}\{(\pm1,0,\pm1,0),\;(0,\pm1,0,\pm1)\}\,.
         \label{eq:C3}
\end{align}
Each $C_i$ contains exactly 8~vertices, and no two subsets share a vertex:
$C_1 \cap C_2 = C_1 \cap C_3 = C_2 \cap C_3 = \varnothing$, with
$C_1 \cup C_2 \cup C_3$ recovering all 24~vertices.

Within each $C_i$, the vertices form a regular 16-cell.  The three 16-cells are
``orthogonal'' in the sense that the inner products between vertices drawn from
different subsets satisfy
\begin{equation}\label{eq:ortho}
  \langle v_i,\, v_j \rangle = \pm\tfrac{1}{\sqrt{2}}
  \quad\text{for}\; v_i \in C_a,\; v_j \in C_b,\; a \neq b\,,
\end{equation}
i.e., vertices from different 16-cells are never parallel or antiparallel.

This triple decomposition is the geometric avatar of \emph{triality}: the three
16-cells are cyclically permuted by an outer automorphism of $D_4$
(Sec.~\ref{sec:triality}).

%======================================================================
\section{The $F_4$ Symmetry Group}\label{sec:F4}
%======================================================================

The automorphism group of the 24-cell is the Weyl group of the exceptional Lie
algebra $\mathfrak{f}_4$:
\begin{equation}\label{eq:WF4}
  \mathrm{Aut}\bigl(\{3,4,3\}\bigr) = W(F_4)\,,
  \qquad |W(F_4)| = 1152\,.
\end{equation}
The order decomposes as $1152 = 2^7 \times 3^2$.  The group acts transitively
on vertices, edges, 2-faces, and 3-cells, reflecting the high degree of
regularity of the 24-cell.

The root system of $F_4$ consists of 48~roots, which are precisely the
24~vertices of the 24-cell together with 24~mid-edge points (forming the dual
24-cell upon rescaling).  In the standard normalisation, the short roots have
length~$1$ and the long roots length~$\sqrt{2}$, with the ratio
$|\alpha_{\text{long}}|/|\alpha_{\text{short}}| = \sqrt{2}$.

The Weyl group $W(F_4)$ contains $W(D_4)$ as a normal subgroup of index~3:
\begin{equation}\label{eq:quotient}
  W(F_4) / W(D_4) \cong S_3\,,
\end{equation}
where $S_3$ is the symmetric group on three elements, acting as the triality
group that permutes the three 16-cell sub-polytopes.  The order of $W(D_4)$ is
$|W(D_4)| = 192 = 2^6 \times 3$, consistent with $1152/192 = 6 = |S_3|$.

%======================================================================
\section{$D_4$ Triality}\label{sec:triality}
%======================================================================

The Lie algebra $\mathfrak{so}(8) \cong D_4$ possesses a unique feature among
all simple Lie algebras: its Dynkin diagram has an $S_3$ symmetry arising from
the threefold branching at the central node.  This symmetry is the \emph{triality
automorphism}, which permutes the three 8-dimensional representations:
\begin{equation}\label{eq:triality}
  \mathbf{8}_v \;\longleftrightarrow\; \mathbf{8}_s
  \;\longleftrightarrow\; \mathbf{8}_c\,,
\end{equation}
where $\mathbf{8}_v$ is the vector representation, $\mathbf{8}_s$ the
positive-chirality spinor, and $\mathbf{8}_c$ the negative-chirality spinor.

In the 24-cell decomposition of Sec.~\ref{sec:decomp}, the three 16-cells
$C_1$, $C_2$, $C_3$ are the geometric realisations of these three
representations.  The triality automorphism acts as
\begin{equation}\label{eq:triality_action}
  \sigma\colon C_1 \to C_2 \to C_3 \to C_1\,,
  \qquad \sigma^3 = \mathrm{id}\,.
\end{equation}

This has a direct physical interpretation in SDGFT: the three 16-cells
correspond to the three fermion generations.  Triality guarantees that the three
generations are treated on an exactly equal footing at the level of the
fundamental polytope, with generation-dependent mass splittings arising only
from the scale-dependent running of the geometric couplings.

The triality automorphism also interchanges the Dynkin labels of $D_4$.
Denoting the four simple roots as $\alpha_1,\alpha_2,\alpha_3,\alpha_4$ with
$\alpha_2$ the central node:
\begin{equation}
  \sigma\colon \alpha_1 \to \alpha_3 \to \alpha_4 \to \alpha_1\,,
  \qquad \alpha_2 \to \alpha_2\,.
\end{equation}

%======================================================================
\section{Dynkin Diagram Structure and the Standard Model}\label{sec:dynkin}
%======================================================================

The $F_4$ Dynkin diagram is
\begin{equation}\label{eq:F4dynkin}
  \underset{\alpha_1}{\circ}\text{---}
  \underset{\alpha_2}{\circ}\text{=\!=}
  \underset{\alpha_3}{\circ}\text{---}
  \underset{\alpha_4}{\circ}
\end{equation}
where the double bond connects $\alpha_2$ and $\alpha_3$ (short--long root
junction).  Removing the peripheral node $\alpha_1$ yields the $B_3$ sub-diagram;
removing $\alpha_4$ yields $C_3$.

The key observation for SDGFT is the chain of maximal-rank sub-algebras
obtained by successive peripheral-node removals and branchings:
\begin{equation}\label{eq:chain}
  F_4 \;\supset\; D_4 \;\supset\;
  A_2 \oplus A_1 \oplus U(1)\,.
\end{equation}
The first step, $F_4 \supset D_4$, corresponds to breaking the long/short root
equivalence.  The second step decomposes $D_4 \cong \mathfrak{so}(8)$ via the
branching rule
\begin{equation}\label{eq:branching}
  \mathfrak{so}(8) \;\supset\;
  \mathfrak{su}(3) \oplus \mathfrak{su}(2) \oplus \mathfrak{u}(1)\,,
\end{equation}
which is precisely the Lie-algebraic content of the Standard Model gauge group
\begin{equation}\label{eq:SM}
  G_{\mathrm{SM}} = SU(3)_C \times SU(2)_L \times U(1)_Y\,.
\end{equation}

The $\mathfrak{su}(3)$ factor (the $A_2$ node) governs the colour interaction,
$\mathfrak{su}(2)$ (the $A_1$ node) the weak isospin, and $\mathfrak{u}(1)$
hypercharge.  No additional $U(1)$ factors or exotic gauge bosons appear: the
rank of $D_4$ is~4, equal to the rank of $G_{\mathrm{SM}}$ when the
hypercharge generator is included.

The representation content follows from the $D_4$ branching rules.  Under the
decomposition~\eqref{eq:branching}, the vector representation splits as
\begin{equation}\label{eq:8v_decomp}
  \mathbf{8}_v \to (\mathbf{3},\mathbf{1})_{+2/3}
               \oplus (\bar{\mathbf{3}},\mathbf{1})_{-2/3}
               \oplus (\mathbf{1},\mathbf{2})_0\,,
\end{equation}
reproducing the quantum numbers of a quark doublet and a lepton doublet.  The
spinor representations $\mathbf{8}_s$ and $\mathbf{8}_c$ decompose analogously,
providing the remaining SM fermion multiplets.

%======================================================================
\section{Conclusions}\label{sec:conclusions}
%======================================================================

We have established the topological and group-theoretic properties of the
24-cell that underpin Scale-Dependent Geometric Field Theory:

\begin{enumerate}
  \item The 24-cell is the unique self-dual regular 4-polytope, with
    $f$-vector $(24,96,96,24)$.

  \item It decomposes into three mutually orthogonal 16-cells
    ($8+8+8=24$), the geometric avatar of $D_4$ triality and the three
    fermion generations.

  \item Its automorphism group $W(F_4)$, of order~1152, contains $W(D_4)$
    as a normal subgroup of index~$6 = |S_3|$.

  \item The Dynkin-diagram chain $F_4 \supset D_4 \supset
    A_2 \oplus A_1 \oplus U(1)$ recovers the Standard Model gauge group
    $SU(3)_C \times SU(2)_L \times U(1)_Y$ without additional assumptions.
\end{enumerate}

Together with the axiomatic foundations of Ref.~\cite{Besemer2026_01} and the
six-cone partition of Ref.~\cite{Besemer2026_03}, these results provide the
geometric basis for the full SDGFT programme.

\begin{thebibliography}{99}

\bibitem{Besemer2026_01}
D.~A.~Besemer,
``Axiomatic Foundations of Scale-Dependent Geometric Field Theory,''
SDGFT-2026-01.

\bibitem{Besemer2026_03}
D.~A.~Besemer,
``Six-Cone Partition of the Three-Sphere and Flavour Topology,''
SDGFT-2026-03.

\bibitem{Coxeter1973}
H.~S.~M.~Coxeter,
\emph{Regular Polytopes}, 3rd ed.\ (Dover, New York, 1973).

\bibitem{Conway1988}
J.~H.~Conway and N.~J.~A.~Sloane,
\emph{Sphere Packings, Lattices and Groups}, 3rd ed.\
(Springer, New York, 1999).

\bibitem{Adams1996}
J.~F.~Adams,
\emph{Lectures on Exceptional Lie Groups}
(University of Chicago Press, Chicago, 1996).

\bibitem{Baez2002}
J.~C.~Baez,
``The octonions,''
Bull.\ Am.\ Math.\ Soc.\ \textbf{39}, 145 (2002).

\bibitem{Humphreys1990}
J.~E.~Humphreys,
\emph{Reflection Groups and Coxeter Groups}
(Cambridge University Press, Cambridge, 1990).

\bibitem{Georgi1974}
H.~Georgi and S.~L.~Glashow,
``Unity of all elementary-particle forces,''
Phys.\ Rev.\ Lett.\ \textbf{32}, 438 (1974).

\end{thebibliography}

\end{document}
